\documentclass{article}

\usepackage{amsmath}

\title{Ch. 27 CW}
\author{Principles of Macroeconomics}
\date{\today}

\begin{document}

\maketitle

\begin{enumerate}

\item What is the likely effect of each of the following events on the stock price of a company? Explain your answers.

	\begin{enumerate}
	
	\item The company announces that while profit is low this year, it has discovered a new line of business that will generate a high profit next year.
	
	\item The company announces that although it had a high profit this year, that profit will be less than had been previously announced.
	
	\item Other companies in the same industry announce that sales are unexpectedly slow this year.
	
	\item The company that announces that it is on track to meet its previously forecast profit target.
	
	\end{enumerate}

\item You have the following information on a closed economy:
	\begin{gather*}
	Y = \$20,000 \\
	C = \$14,000 \\
	G = \$2,500 \\
	T = \$1,700
	\end{gather*}
	Calculate private saving, public saving, national saving, and investment.
	
\item A firm is considering an investment in a machine that will enable the firm to earn \$500 in revenue next year and each subsequent year. The machine is super durable and does not depreciate at all. The firm's next best alternative is to save the money at a 10\% annual interest rate. The inflation rate is zero.

	\begin{enumerate}
	
	\item If the price of the machine is \$4,500 today, should the firm make the investment in the machine or not? Briefly explain why or why not.
	
	\item Now suppose that before the firm makes the decision, the prevailing interest rate on savings accounts increases to 12\%, but there is no change in the expected revenue from the investment. Does the machine remain a worthwhile investment?
	
	\item What is the maximum price the firm should be willing to pay for the machine if the interest rate on savings is 10\% (in other words, at what price is the firm indifferent between the two options)?
	
	\item Does a higher interest rate encourage or discourage investment? Does it encourage or discourage saving? 
	
	\end{enumerate}
	
\item For each of the following scenarios, predict the resulting changes in the real interest rate, national saving, and investment. Show all your diagrams, and write a two sentence explanation for why there's a shift and how the shift impacts the market.

	\begin{enumerate}
	
	\item A reduction in military spending moves the government’s budget from deficit into surplus.
	
	\item The government raises its tax on corporate profits. Other tax changes are also made such that the government’s deficit remains unchanged.
	
	\item Consumer's become optimistic about the economy and decrease their precautionary saving. 
	
	\item A new generation of computer-controlled machines becomes available. These machines produce manufactured goods much more quickly and with fewer defects.
	
	\item New environmental regulations increase firms’ costs of operating capital.
	
	\item The government decreases its tax on corporate profits, but it doesn't cut government spending.
	
	\end{enumerate}

\end{enumerate}

\end{document}